% !PRINTED 502
\weberpage{502}
%<<CONT

3. \gesperrt{Der absolute Werth eines Productes zweier oder
mehrerer rationaler Functionale ist gleich dem Producte der
absoluten Werthe der Factoren.}

Ist $\omega$ ein Functional des Körpers $\Omega$, so ist $N(\omega)$ ein rationales
Functional, und die \gesperrt{Norm eines Productes oder eines Quotienten
zweier Functionale ist gleich dem Producte oder dem Quotienten
der Normen der Bestandtheile.}

4. \gesperrt{Wir nennen den absoluten Werth des rationalen
Functionales $N(\omega)$ die absolute Norm $N_a(\omega)$ von
$\omega$ und setzen}
\weq{3}{N(\omega) = N_a(\omega)\,E(\omega).}
$N_a(\omega)$ ist immer eine positive rationale Zahl, und $E(\omega)$
ist der Quotient zweier primitiver ganzer Functionen. Aus 3.
folgt, wenn $\alpha, \beta$ zwei Functionale in $\Omega$ sind, die Formel
\weq{4}{N_a(\alpha\beta) = N_a(\alpha)\,N_a(\beta),}
oder der Satz:

5. \gesperrt{Die absolute Norm eines Productes von Functionalen
in $\Omega$ ist gleich dem Producte der absoluten Normen der
Factoren.}

Wenn wir alle Co\"efficienten eines Functionales in $\Omega$ durch
die entsprechenden Zahlen eines zu $\Omega$ conjugirten Körpers $\Omega_1$
ersetzen, so erhalten wir ein \gesperrt{conjugirtes Functional}. Der
gesammte Functionalkörper $\overline{\Omega}$ geht dadurch in einen
\gesperrt{conjugirten Functionalkörper} $\overline{\Omega}_1$ über. Da jede Gleichung
zwischen Functionalen des Körpers $\Omega$ im Grunde nur die
Zusammenfassung einer Reihe von Gleichungen zwischen Zahlen des
Körpers $\Omega$ ist, so haben wir den Satz (Bd. I, \S.\,141):

6. \gesperrt{Jede Gleichung zwischen Functionalen des Körpers
$\overline{\Omega}$ bleibt richtig, wenn für alle Functionale die
entsprechenden Elemente eines conjugirten Körpers $\overline{\Omega}_1$
gesetzt werden.}

Bedeutet $t$ eine in $\omega$ nicht vorkommende Variable, so hat
die Function $N(t-\omega)$ die Form
\weq{5}{\Phi(t) = t^n + A_1 t^{n-1} + A_2 t^{n-2} + \cdots + A_n,}
worin die Co\"efficienten $A_i$ rationale Functionale sind, und diese
Function $\Phi(t)$ verschwindet, wenn $\omega$ für $t$ gesetzt wird. Es giebt
aber nicht bloss eine solche Function $\Phi(t)$, die für $t = \omega$ ver-
%CONT>>
